\section{Introduction}
\label{sec:introduction}

Express.js is among the most widely used web frameworks for Node.js, and almost every non-trivial
service built on it reaches a relational database through some \emph{access layer}. Access layers
trade raw performance for ergonomics, type safety, and protection from pitfalls such as the N+1 query
problem~\cite{chen2014orm} (Section~\ref{sec:related}). Native drivers expose the wire protocol
directly (\texttt{pg}, \texttt{mysql2}), query builders assemble SQL programmatically (Knex), and
object-relational mappers (ORMs) map rows to objects following patterns such as Data Mapper and
Active Record~\cite{fowler2002patterns}. Node.js back-end performance under load is a recognized
empirical topic~\cite{chaniotis2015nodejs,kuffel2025nodejs,carmo2024backend}, but that literature
varies the runtime and framework, not the access layer.

The most visible comparisons are vendor-authored, and the peer-reviewed ones surveyed in
Section~\ref{sec:related} each leave part of the space uncovered: breadth of products, both engines,
a reported tail, or measurement through an HTTP framework (Supplement Table~S51).
These gaps motivate the coverage of this experiment, not a priority claim.
A more basic gap is methodological: the reported numbers are not established to be
\emph{comparable}, since a layer returning fewer fields or erroring faster than it works can look
faster, a library's chosen query strategy is conflated with its own cost, and a single, arbitrary
load point conflates capacity, tail latency under demand, and per-request latency.
Establishing that comparability is the problem this paper takes up, and it is a
software-engineering one: where such benchmarks inform a technology choice, a comparison not
established to be comparable can misdirect that choice rather than merely misreport a number.

This paper addresses these gaps with a \emph{comparability protocol} (a structured, executable
proposal rather than a validated methodology) and one non-vendor-authored empirical instantiation of
it: an \emph{identical} Express application served through all nine access layers over five selected
access patterns, which stress distinct aspects of the access layer and are not a representative
sample of Express and database workloads. Seven layers run against \emph{both} engines under an
identical schema, dataset, and pool size; two tuned native-driver baselines are disclosed reference
points, not an optimization ceiling. Every (layer, engine, pattern) reports throughput and the
high-load p99 under equal closed-loop demand at a fixed 50-connection point over 25 runs. That point
sits near saturation, where the p99 largely tracks capacity and queueing, so we also report the tail
at \emph{matched utilization}.
The study separates two estimands: the \emph{policy-selected documented configuration}
performance of each layer (RQ1--RQ3), and a \emph{common-SQL raw-path sensitivity} contrast reporting
the residual under identical SQL, which changes query formulation, protocol, preparation, round trips,
and mapping together rather than isolating any one. The full harness (adapters, seed data, an
autocannon-driven~\cite{autocannon} load runner, a specification-derived expected-result oracle,
differential-equivalence checks, and write-state validation) is released as an open replication
package.\footnote{Harness, deterministic seed, all adapters, raw per-cell measurements, and the
table/figure-generating scripts: \texttt{https://github.com/mmiotk/express-db-access-performance}.
The archived release and its DOIs are given in the Data availability statement.}

The experiment answers three research questions, each scoped to the benchmarked
implementations under the specified operating point:
\begin{enumerate}[label=\textbf{RQ\arabic*},leftmargin=*]
  \item \emph{(performance difference):} How large is the throughput and response-time-p99
        \emph{difference} relative to the native-driver baseline, and how does it vary across the
        five access patterns? We keep distinct three quantities a fixed operating point conflates:
        \emph{capacity}, \emph{high-load latency}, and \emph{latency at matched utilization}.
  \item \emph{(backend-stack transfer):} Does the observed layer ordering transfer between the
        tested PostgreSQL and MySQL backend stacks, where changing the DBMS can also change its
        supported adapter and lower-level driver? The question concerns transfer, not causation: the
        design cannot attribute a failure to transfer to the DBMS rather than to the adapter or
        driver the backend brings with it. We answer it at two levels, whether a specific pair
        reverses durably under a stated rule and whether the whole ordering agrees, the latter
        reported descriptively.
  \item \emph{(pattern spread):} Which of the five tested patterns shows the
        largest native-relative spread in this configuration?
\end{enumerate}
Supplement Table~S47 states, per question, the quantity the design estimates and the quantity it
does \emph{not} identify.

This paper's primary contribution is a \textbf{comparability protocol for access-layer
benchmarking}, specified normatively in Section~\ref{sec:protocol} (inputs, stages, cell-admission
criteria, interpretation, and limits).
It rests on three constructs, defined there and named here.
\emph{Semantic admission} requires evidence that a treatment performed the declared task before any
of its timings may be used, so a layer that returns less, writes wrong, or errors faster than it
works cannot appear faster.
\emph{Treatment definition} fixes each treatment by a predeclared, reproducible selection rule,
which prevents result-dependent tuning but establishes nothing about whether the configuration is
optimal, common, or preferred in production.
\emph{Operating-point separation} reports capacity, high-load tail under equal demand, and tail at
matched utilization as \emph{distinct} quantities.
These constituent controls are established performance-evaluation and testing
ideas~\cite{raasveldt2018fair,fruth2022tail,georges2007rigorous}, and CYNTHIA already applies
differential equivalence across ORM implementations~\cite{sotiropoulos2021cynthia}; the contribution
is their integration and operationalization for relational access-layer comparison.
We \emph{exercise} that discipline through an in-experiment ablation and a single-rater audit of
nine external benchmarks; neither is independent validation, and Section~\ref{sec:discussion} sets
out what the exercise does and does not establish.

Two contributions should be kept apart. The \emph{methodological} contribution is the protocol
and its checklist: proposals for how access-layer comparisons should be conducted, reported,
and audited. The \emph{empirical} contribution is one instantiation of them, scoped to the
conditions and software versions of Section~\ref{sec:methodology}. No result is offered as a
property of a library, and the study measures no non-performance quality such as maintainability,
type safety, migration tooling, or ecosystem maturity.
